\documentclass{article}

\usepackage{tabularx}
\usepackage{booktabs}

\title{Problem Statement and Goals\\\progname}

\author{\authname}

\date{}

\input{../Comments.text}
\input{../Common.text}

\begin{document}

\maketitle

\begin{table}[hp]
\caption{Revision History} \label{TblRevisionHistory}
\begin{tabularx}{\textwidth}{llX}
\toprule
\textbf{Date} & \textbf{Developer(s)} & \textbf{Change}\\
\midrule
Spetember 22 & Justin Kwinecki & Initial Problem Statement and Goals\\
Date2 & Name(s) & Description of changes\\
... & ... & ...\\
\bottomrule
\end{tabularx}
\end{table}

\section{Problem Statement}

\wss{You should check your problem statement with the
\href{https://smiths.github.io/capTemplate/Checklists/ProbState-Checklist.pdf}
{problem statement checklist}}. 

\wss{You can change the section headings, as long as you include the required
information.}

\subsection{Problem}

Modern LLM applications increasingly rely on multiple agents that interact with the same underlying context.
These agents often process a shared prompt or prompt tree, causing the same prefix information to be recomputed 
independently for each request. As the number of simultaneous agents increases, this redundant computation 
increases inference latency, wastes computational resources, and limits the scalability of multi-agent LLM systems.

The problem is to extend an LLM serving engine so that requests from multiple agents can efficiently recognize and 
exploit shared prompt prefixes during prefill. The system should combine shared portions of the prompt tree into a 
single prefill operation and coordinate simultaneous requests so that shared prefixes are computed once rather than 
repeatedly.

\subsection{Inputs and Outputs}

\wss{Characterize the problem in terms of ``high level'' inputs and outputs.  
Use abstraction so that you can avoid details.}

\subsection{Stakeholders}

\subsection{Environment}

\wss{Hardware and Software Environment for the users.  The developer environment
is summarized as part of the developer plan.}

\section{Goals}

\begin{itemize}

\item \textbf{Reduce redundant computation:} Reduce the amount of computation required to prefill shared prompt context across simultaneous agent requests.
\item \textbf{Reduce inference latency:} Decrease prefill latency for multi-agent workloads with overlapping context.
\item \textbf{Improve serving throughput:} Increase the number of multi-agent requests that can be served within a given period when requests share prompt context.
\item \textbf{Preserve correctness: Ensure:} that sharing prefill computation does not change the resulting model behavior compared with independent request processing.
\item \textbf{Maintain existing functionality:} Continue to support requests without shared context through the existing serving workflow.

\end{itemize}

\section{Stretch Goals}

\begin{itemize}
\item \textbf{Improve scalability:} Demonstrate that the benefits increase as the number of simultaneous agents sharing context grows.
\item \textbf{Support dynamic prompt trees:} Efficiently handle agents that branch from shared context at different points in their interactions.
\item \textbf{Optimize resource utilization:} Reduce unnecessary GPU computation or memory usage resulting from repeated processing of shared context.
\item \textbf{Evaluate broader workloads:} Measure performance across different prompt sizes, degrees of prefix overlap, and numbers of concurrent agents.
\item \textbf{Demonstrate substantial performance gains:} Establish measurable improvements over independent prefill across representative multi-agent workloads.
\end{itemize}

\section{Custom Documents (CDs)}

\wss{For CAS 741: State whether the project is a research project. This
designation, with the approval (or request) of the instructor, can be modified
over the course of the term.}

\wss{For SE Capstone: List your custom documents (CDs).  Potential CDs include
usability testing, code walkthroughs, user documentation, formal proof,
GenderMag personas, Design Thinking, etc.  (The full list is on the course
outline and in Lecture 02.) Normally the number of CDs will be two (one per
term).  Approval of the CDs will be part of the discussion with the instructor
for approving the project. The CDs, with the approval (or request) of the
instructor, can be modified over the course of the term.}

\newpage{}

\section*{Appendix --- Reflection}

\wss{Not required for CAS 741}

\input{../Reflection.text}

\begin{enumerate}
    \item What went well while writing this deliverable? 
    \item What pain points did you experience during this deliverable, and how
    did you resolve them?
    \item How did you and your team adjust the scope of your goals to ensure
    they are suitable for a Capstone project (not overly ambitious but also of
    appropriate complexity for a senior design project)?
\end{enumerate}  

\end{document}